% !TEX root = main.tex
\chapter[Extreme Amenability and Ramsey Theory]{Extreme Amenability and the Ramsey Property}
\label{chap:extreme-amenability-ramsey}

The aim of this chapter is to connect the dynamical notion of extreme amenability introduced in the previous chapter with the Ramsey property for classes of finite structures. This is the first decisive step in the Kechris--Pestov--Todorcevic correspondence. For closed subgroups of \(S_\infty\), the topology is controlled by stabilisers of finite sets, and this makes it possible to reformulate a fixed-point property for flows in terms of finite colourings of finite configurations. The result is a bridge from topological dynamics to structural Ramsey theory.

We begin by recalling the basic notions of structural Ramsey theory and some standard examples. We then prove a characterisation of extremely amenable closed subgroups of \(S_\infty\) in terms of a Ramsey property for the corresponding finite types. Finally, we translate this characterisation into Fraïssé-theoretic language and obtain the fundamental criterion that if \(\bff\) is the Fraïssé limit of a Fraïssé order class \(\cK\), then \(\Aut(\bff)\) is extremely amenable if and only if \(\cK\) has the Ramsey property. This criterion will be the main tool in the next chapter, where we construct concrete examples of extremely amenable automorphism groups.


\section{Structural Ramsey theory}
\label{sec:structural-ramsey-theory}

We begin by recalling the basic notions of structural Ramsey theory that will be used in the sequel. The point of this section is to formulate the combinatorial property that later characterises extreme amenability for automorphism groups of Fraïssé limits.

Let \(\bfa\) and \(\bfb\) be \(L\)-structures. If \(\bfa\hookrightarrow \bfb\), we let
\[
    \begin{pmatrix} \bfb \\ \bfa \end{pmatrix}
    =
    \{\bfa' \leq \bfb : \bfa' \cong \bfa\},
\]
the set of all substructures of \(\bfb\) isomorphic to \(\bfa\).

If \(\bfa\hookrightarrow \bfb\hookrightarrow \bfc\) and \(k \geq 2\), we write \(\bfc \to (\bfb)^{\bfa}_k\) to mean that for every colouring
\(c \colon \begin{pmatrix} \bfc \\ \bfa \end{pmatrix} \to \{1,\dots,k\}\), there is some
\(\bfb' \in \begin{pmatrix} \bfc \\ \bfb \end{pmatrix}\) such that
\(c \restriction \begin{pmatrix} \bfb' \\ \bfa \end{pmatrix}\) is constant. In this situation we say that \(\bfb'\) is \emph{homogeneous} for the colouring \(c\).

\begin{definition}
    Let \(\cK\) be a class of finite \(L\)-structures. We say that \(\cK\) has the \emph{Ramsey property} if \(\cK\) is hereditary and for every \(\bfa,\bfb \in \cK\) with \(\bfa\hookrightarrow \bfb\) and every integer \(k \geq 2\), there exists \(\bfc \in \cK\) such that
    \[
        \bfc \to (\bfb)^{\bfa}_k.
    \]
\end{definition}

As usual, it is enough to verify the Ramsey property for two colours.

\begin{proposition}
    A hereditary class \(\cK\) of finite structures has the Ramsey property if and only if for every \(\bfa,\bfb \in \cK\) with \(\bfa\hookrightarrow \bfb\), there exists \(\bfc \in \cK\) such that
    \[
        \bfc \to (\bfb)^{\bfa}_2.
    \]
\end{proposition}

\begin{proof}
    Only one implication needs proof. Assume the two-colour condition holds. We show by induction on \(k \geq 2\) that for all \(\bfa\hookrightarrow \bfb\) in \(\cK\) there exists \(\bfc \in \cK\) with \(\bfc \to (\bfb)^{\bfa}_k\).

    The case \(k=2\) is our assumption. Suppose the statement has been proved for \(k-1\), where \(k \geq 3\), and fix \(\bfa\hookrightarrow \bfb\) in \(\cK\). Choose \(\bfd \in \cK\) such that \(\bfd \to (\bfb)^{\bfa}_{k-1}\), and then choose \(\bfc \in \cK\) such that \(\bfc \to (\bfd)^{\bfa}_2\).

    Now let \(c \colon \begin{pmatrix} \bfc \\ \bfa \end{pmatrix} \to \{1,\dots,k\}\) be a \(k\)-colouring. Define a two-colouring \(c' \colon \begin{pmatrix} \bfc \\ \bfa \end{pmatrix} \to \{0,1\}\) by
    \[
        c'(\bfa')=
        \begin{cases}
            0, & c(\bfa')=k,                 \\
            1, & c(\bfa')\in\{1,\dots,k-1\}.
        \end{cases}
    \]
    Since \(\bfc \to (\bfd)^{\bfa}_2\), there is \(\bfd' \in \begin{pmatrix} \bfc \\ \bfd \end{pmatrix}\) homogeneous for \(c'\). If \(c'\) is constantly \(0\) on \(\begin{pmatrix} \bfd' \\ \bfa \end{pmatrix}\), then every copy of \(\bfa\) in \(\bfd'\) has colour \(k\), so \(\bfd'\) contains a homogeneous copy of \(\bfb\). If \(c'\) is constantly \(1\), then the original colouring \(c\) uses only the colours \(1,\dots,k-1\) on \(\begin{pmatrix} \bfd' \\ \bfa \end{pmatrix}\). Since \(\bfd' \cong \bfd\) and \(\bfd \to (\bfb)^{\bfa}_{k-1}\), there is a copy of \(\bfb\) inside \(\bfd'\) homogeneous for \(c\).

    Thus \(\bfc \to (\bfb)^{\bfa}_k\), as required.
\end{proof}

We now record several standard examples.

\begin{example}
    The class \(\cLO\) of finite linear orders has the Ramsey property.
\end{example}

\begin{proof}
    Two finite linear orders are isomorphic if and only if they have the same cardinality. Thus, for finite linear orders \(\bfa\), \(\bfb\), and \(\bfc\), the set
    \(\begin{pmatrix} \bfc \\ \bfa \end{pmatrix}\) may be identified with the family of \(|A|\)-element subsets of \(C\). Under this identification, the statement \(\bfc \to (\bfb)^{\bfa}_k\) is exactly the classical finite Ramsey theorem.
\end{proof}

\begin{example}
    Let \(L=\{<,E\}\), where \(<\) is a binary relation symbol and \(E\) is a binary relation symbol interpreted as a symmetric irreflexive graph relation. Then the class \(\vec{\mathcal{GR}}\) of finite ordered graphs has the Ramsey property.
\end{example}

This is the theorem of Ne\v{s}et\v{r}il and R\"odl~\cite{NR77,NR83}.

\begin{example}
    Let \(F\) be a finite field, and let \(\mathcal{V}_F\) be the class of finite-dimensional vector spaces over \(F\), regarded as structures in a language coding addition and scalar multiplication. Then \(\mathcal{V}_F\) has the Ramsey property.
\end{example}

This is the Graham--Leeb--Rothschild theorem~\cite{GLR}.

\begin{example}
    Let \(\mathcal{BA}\) be the class of finite Boolean algebras in the usual algebraic language. Then \(\mathcal{BA}\) has the Ramsey property.
\end{example}

This is an equivalent formulation of the Dual Ramsey Theorem of Graham and Rothschild; see, for example,~\cite[\S 4.13]{N95}.

\begin{definition}
    A structure \(\bfa\) is said to be \emph{rigid} if every automorphism of \(\bfa\) is the identity. Equivalently, \(\Aut(\bfa)=\{\id_{\bfa}\}\).
    A class of finite structures is called rigid if every member of the class is rigid.
\end{definition}

For our purposes, ordered classes are especially important, since finite linearly ordered structures are rigid. This rigidity ties the Ramsey property to amalgamation.

\begin{proposition}
    \label{prop:rigid-rp-implies-ap}
    Let \(\cK\) be a hereditary class of finite rigid \(L\)-structures. If \(\cK\) has the joint embedding property and the Ramsey property, then \(\cK\) has the amalgamation property.
\end{proposition}

\begin{proof}
    Let \(\bfa,\bfb,\bfc \in \cK\) and let \(f \colon \bfa \hookrightarrow \bfb\) and \(g \colon \bfa \hookrightarrow \bfc\) be embeddings. Since \(\cK\) has the joint embedding property, there exists \(\bfe \in \cK\) into which both \(\bfb\) and \(\bfc\) embed. Fix embeddings \(u \colon \bfb \hookrightarrow \bfe\) and \(v \colon \bfc \hookrightarrow \bfe\). By the Ramsey property, there is \(\bfd \in \cK\) such that \(\bfd \to (\bfe)^{\bfa}_4\).

    We define a colouring of \(\begin{pmatrix} \bfd \\ \bfa \end{pmatrix}\) with values in the four-element set \(\mathcal{P}(\{\bfb,\bfc\})\). For \(\bfa' \in \begin{pmatrix} \bfd \\ \bfa \end{pmatrix}\), put
    \[
        \bfb \in c(\bfa')
        \quad\Longleftrightarrow\quad
        \text{there exists an embedding } r \colon \bfb \hookrightarrow \bfd \text{ with } r(f(A))=A',
    \]
    and define similarly
    \[
        \bfc \in c(\bfa')
        \quad\Longleftrightarrow\quad
        \text{there exists an embedding } s \colon \bfc \hookrightarrow \bfd \text{ with } s(g(A))=A'.
    \]

    Choose \(\bfe' \in \begin{pmatrix} \bfd \\ \bfe \end{pmatrix}\) homogeneous for this colouring. Since \(\bfe' \cong \bfe\) contains copies of both \(\bfb\) and \(\bfc\), there is some copy of \(\bfa\) inside \(\bfe'\) whose colour contains \(\bfb\), and also some copy of \(\bfa\) inside \(\bfe'\) whose colour contains \(\bfc\). Because \(\bfe'\) is homogeneous, the common colour of all members of \(\begin{pmatrix} \bfe' \\ \bfa \end{pmatrix}\) must therefore be exactly \(\{\bfb,\bfc\}\).

    Now choose any \(\bfa' \in \begin{pmatrix} \bfe' \\ \bfa \end{pmatrix}\). Since \(c(\bfa')=\{\bfb,\bfc\}\), there are embeddings \(r \colon \bfb \hookrightarrow \bfd\) and \(s \colon \bfc \hookrightarrow \bfd\) such that \(r(f(A))=A'=s(g(A))\). Hence both \(r \circ f\) and \(s \circ g\) are isomorphisms from \(\bfa\) onto \(\bfa'\). Since \(\bfa\) is rigid, this isomorphism is unique, so \(r \circ f = s \circ g\).
    Thus \(\bfd\), together with the embeddings \(r\) and \(s\), is an amalgam of \(\bfb\) and \(\bfc\) over \(\bfa\). Therefore \(\cK\) has the amalgamation property.
\end{proof}

In particular, if \(\cK\) is a nonempty countable hereditary class of finite rigid structures, has the joint embedding property, contains structures of arbitrarily large cardinality, and has the Ramsey property, then \(\cK\) is a Fraïssé class.

The importance of these observations for us is that many of the classes arising in the Kechris--Pestov--Todorcevic correspondence are order classes. For such classes, rigidity is automatic, and Proposition~\ref{prop:rigid-rp-implies-ap} shows that once the Ramsey property is established, the amalgamation property follows under mild additional hypotheses.



\section{Extreme amenability for closed subgroups of \texorpdfstring{\(S_\infty\)}{S-infinity}}

\label{sec:ea-closed-subgroups-sinf}

We now reformulate extreme amenability for closed subgroups of \(S_\infty\) in combinatorial terms. The basic idea is that for such groups, open subgroups are determined by finite information, and this allows one to express a fixed-point property for flows as a Ramsey-type statement about finite configurations.

We begin with a general criterion for the existence of a fixed point.

\begin{lemma}
    \label{lem:fixed-point-criterion}
    Let \(G\) be a topological group and let \(X\) be a \(G\)-flow. Then the following are equivalent.
    \begin{enumerate}
        \item[(i)] \(X\) has a fixed point.
        \item[(ii)] For every \(n \geq 1\), every continuous map \(f \colon X \to \R^n\), every \(\epsilon >0\), and every finite set \(F \subseteq G\), there exists \(x \in X\) such that
              \[
                  \norm{f(x)-f(g\cdot x)}\leq \epsilon
                  \qquad\text{for all } g\in F,
              \]
              where \(\norm{\cdot}\) denotes the Euclidean norm.
    \end{enumerate}
\end{lemma}

\begin{proof}
    The implication (i) \(\Rightarrow\) (ii) is immediate.

    For (ii) \(\Rightarrow\) (i), define, for each continuous \(f \colon X \to \R^n\), each \(\epsilon>0\), and each finite \(F\subseteq G\),
    \[
        A_{f,\epsilon,F}
        =
        \{x\in X : \norm{f(x)-f(g\cdot x)}\leq \epsilon \text{ for all } g\in F\}.
    \]
    Each set \(A_{f,\epsilon,F}\) is closed.
    We claim that the family \(\{A_{f,\epsilon,F}\}\) has the finite intersection property.
    Indeed, given finitely many triples \((f_j,\epsilon_j,F_j)\), \(j=1,\dots,m\), let \(\bar F=F_1\cup\cdots\cup F_m\) and \(\bar\epsilon=\min\{\epsilon_1,\dots,\epsilon_m\}\), and define \(\bar f=(f_1,\dots,f_m)\colon X \to \R^{n_1+\cdots+n_m}\), where \(f_j\colon X\to\R^{n_j}\).
    Then
    \[
        A_{\bar f,\bar\epsilon,\bar F}
        \subseteq
        \bigcap_{j=1}^m A_{f_j,\epsilon_j,F_j},
    \]
    and \(A_{\bar f,\bar\epsilon,\bar F}\) is nonempty by (ii).
    Since \(X\) is compact, it follows that
    \[
        \bigcap_{f,\epsilon,F} A_{f,\epsilon,F}\neq\emptyset.
    \]
    Choose
    \[
        x\in \bigcap_{f,\epsilon,F} A_{f,\epsilon,F}.
    \]
    We show that \(x\) is fixed. Suppose not. Then there exists \(g\in G\) such that \(g\cdot x\neq x\). By complete regularity of the compact Hausdorff space \(X\), there is a continuous function \(f\colon X\to\R\) such that \(f(x)=0\) and \(f(g\cdot x)=1\).
    But then \(x\notin A_{f,1/2,\{g\}}\), a contradiction. Hence \(x\) is a fixed point.
\end{proof}

We now specialize to closed subgroups of \(S_\infty\).

\begin{proposition}
    \label{prop:coset-characterization}
    Let \(G\leq S_\infty\) be a closed subgroup. Then the following are equivalent.
    \begin{enumerate}
        \item[(i)] \(G\) is extremely amenable.
        \item[(ii)] For every open subgroup \(V\leq G\), every finite colouring \(c\colon G/V\to\{1,\dots,k\}\), and every finite set \(A\subseteq G/V\), there exist \(g\in G\) and \(i\in\{1,\dots,k\}\) such that
              \[
                  c(g\cdot a)=i
                  \qquad\text{for all } a\in A,
              \]
              where \(G\) acts on \(G/V\) by left translation.
    \end{enumerate}
\end{proposition}

\begin{proof}
    Assume first that \(G\) is extremely amenable. Let \(V\leq G\) be open, let \(c\colon G/V\to\{1,\dots,k\}\) be a colouring, and let \(A\subseteq G/V\) be finite. Consider the compact space \(Y=\{1,\dots,k\}^{G/V}\) with the product topology, endowed with the shift action
    \[
        (g\cdot p)(x)=p(g^{-1}\cdot x)
        \qquad
        (g\in G,\ p\in Y,\ x\in G/V).
    \]
    Let \(X=\overline{G\cdot c}\subseteq Y\). Since \(G\) is extremely amenable, \(X\) has a fixed point \(\gamma\). Because the action of \(G\) on \(G/V\) is transitive,
    every fixed point of \(Y\) is constant.
    Thus there is \(i\in\{1,\dots,k\}\) such that \(\gamma(x)=i\) for all \(x\in G/V\). Since \(\gamma\in\overline{G\cdot c}\), there is \(g\in G\) such that \((g^{-1}\cdot c)(a)=\gamma(a)=i\) for all \(a\in A\).
    Equivalently, \(c(g\cdot a)=i\) for all \(a\in A\).

    Conversely, assume (ii), and let \(X\) be a \(G\)-flow. By Lemma~\ref{lem:fixed-point-criterion}, it is enough to verify condition (ii) of that lemma. So fix a continuous map \(f\colon X\to\R^n\), an \(\epsilon>0\), and a finite set \(F\subseteq G\).

    By continuity of the action and compactness of \(X\), there exists an open neighbourhood \(U\) of the identity in \(G\) such that
    \[
        \norm{f(x)-f(u\cdot x)}\leq \epsilon/3
        \qquad\text{for all } x\in X,\ u\in U.
    \]
    Since \(G\) is a subgroup of \(S_\infty\) equipped with the subspace topology, it has a neighbourhood basis at the identity consisting of open subgroups.
    Choose an open subgroup \(V\leq U\).

    Partition the compact set \(f(X)\subseteq\R^n\) into finitely many sets \(A_1,\dots,A_k\) of diameter at most \(\epsilon/3\). Fix \(x_0\in X\) and define \(U_i=\{g\in G : f(g\cdot x_0)\in A_i\}\). Let \(W_i=V U_i=\{vg : v\in V,\ g\in U_i\}\).
    Each \(W_i\) is a union of right cosets of \(V\). Equivalently, each \(W_i^{-1}\) is a union of left cosets of \(V\). Since \(G=U_1\cup\cdots\cup U_k\), we also have \(G=W_1\cup\cdots\cup W_k\).
    Define a colouring \(c\colon G/V\to\{1,\dots,k\}\) so that
    \[
        c(gV)=i
        \quad\Longrightarrow\quad
        g^{-1}\in W_i.
    \]
    This is possible because the sets \(W_i^{-1}\) cover \(G\) and each is a union of left cosets of \(V\).

    Apply (ii) to the finite subset \(A'=\{hV : h\in F\cup\{1_G\}\}\subseteq G/V\). There exist \(g\in G\) and \(i\in\{1,\dots,k\}\) such that \(c(ghV)=i\) for all \(h\in F\cup\{1_G\}\). By the definition of \(c\), this means that \(h^{-1}g^{-1}\in W_i=V U_i\) for all \(h\in F\cup\{1_G\}\), or equivalently, \(hg\in U_i^{-1}V\).
    Replacing \(U_i\) by \(U_i^{-1}\) in the construction if necessary, we may assume that for each \(h\in F\cup\{1_G\}\) there exists \(v_h\in V\) such that \(v_h h g \in U_i\). Set \(x=g\cdot x_0\). Then for each \(h\in F\cup\{1_G\}\) we have \(f(v_h h\cdot x)=f(v_h h g\cdot x_0)\in A_i\). Since \(v_h\in V\subseteq U\), it follows that \(\norm{f(h\cdot x)-f(v_h h\cdot x)}\leq \epsilon/3\).
    Hence \(f(h\cdot x)\) lies in the \(\epsilon/3\)-neighbourhood of \(A_i\) for every \(h\in F\cup\{1_G\}\). As the diameter of \(A_i\) is at most \(\epsilon/3\), we obtain
    \[
        \norm{f(x)-f(h\cdot x)}\leq \epsilon
        \qquad\text{for all } h\in F.
    \]
    Thus \(X\) has a fixed point, and \(G\) is extremely amenable.
\end{proof}

For closed subgroups of \(S_\infty\), it is enough to consider pointwise stabilisers of finite sets. For each nonempty finite \(F\subseteq \N\), let
\[
    G_{(F)}=\{g\in G : g(n)=n \text{ for all } n\in F\}.
\]
These subgroups form a neighbourhood basis at the identity.

We now recast the preceding proposition in terms of finite subsets of \(\N\). The action of \(G\) on \(\N\) induces an action on the set of all finite subsets of \(\N\) by \(g\cdot F=\{g(n): n\in F\}\).
For a finite nonempty set \(F\subseteq\N\), let
\[
    G_{\{F\}}=\{g\in G : g\cdot F=F\}
\]
be its setwise stabiliser. The orbit \(G\cdot F\) will be called the \emph{\(G\)-type} of \(F\). If \(\rho\) and \(\sigma\) are \(G\)-types, we write \(\rho\leq \sigma\) if some, equivalently every, member of \(\sigma\) contains some member of \(\rho\).

We also need the logic action on the space of structures on \(\N\). If \(L\) is a countable signature, let \(X_L\) denote the space of all \(L\)-structures with universe \(\N\). When \(L\) is relational, \(X_L\) is compact, and the natural action of \(S_\infty\) on \(X_L\) is continuous.

\begin{draftnote}
    Suppose that \(L=\{R_i:i\in I\}\) is relational, where \(R_i\) has arity
    \(n_i\). An \(L\)-structure \(\bfa\) with universe \(\N\) is completely
    determined by the truth values
    \[
        R_i^{\bfa}(\bar a),
        \qquad
        i\in I,\quad \bar a\in\N^{n_i}.
    \]
    Hence
    \[
        X_L
        \cong
        \prod_{i\in I}\{0,1\}^{\N^{n_i}},
    \]
    where \(\{0,1\}\) carries the discrete topology. A basic open set specifies
    the truth values of only finitely many atomic statements. Since
    \(\{0,1\}\) is compact, Tychonoff's theorem implies that \(X_L\) is compact.

    The logic action of \(S_\infty\) on \(X_L\) is given by
    \[
        R_i^{g\cdot\bfa}(\bar a)
        \quad\Longleftrightarrow\quad
        R_i^{\bfa}(g^{-1}\bar a).
    \]
    To see that this action is continuous, fix \(g\in S_\infty\),
    \(\bfa\in X_L\), and a basic neighbourhood of \(g\cdot\bfa\), determined
    by finitely many atomic statements \(R_i(\bar a)\). Let \(A\subseteq\N\)
    be the finite set of all entries occurring in these tuples. Require
    \(h\in S_\infty\) to agree with \(g\) on \(g^{-1}(A)\), and require
    \(\bfb\in X_L\) to agree with \(\bfa\) on the corresponding finitely many
    statements \(R_i(g^{-1}\bar a)\). Then
    \(h^{-1}\bar a=g^{-1}\bar a\), and therefore
    \[
        R_i^{h\cdot\bfb}(\bar a)
        \iff
        R_i^{\bfb}(g^{-1}\bar a)
        \iff
        R_i^{\bfa}(g^{-1}\bar a)
        \iff
        R_i^{g\cdot\bfa}(\bar a).
    \]
    Thus a neighbourhood of \((g,\bfa)\) is mapped into the prescribed
    neighbourhood of \(g\cdot\bfa\), proving joint continuity.

    The relational hypothesis is essential for compactness. A function symbol
    of arity \(n\) would be interpreted by an element of
    \(\N^{\N^n}\), and a constant symbol by an element of the infinite discrete
    space \(\N\); these spaces are not compact. In particular,
    \(\LO(\N)\) is compact because it is a closed subspace of
    \(X_{\{<\}}\): every failure of the axioms of a linear order is witnessed
    on finitely many elements.
\end{draftnote}

In the special case \(L=\{<\}\), let \(\LO(\N)\subseteq X_L\) be the compact space of all linear orderings on \(\N\). We say that \(G\) \emph{preserves an ordering} if the \(G\)-flow \(\LO(\N)\) has a fixed point, that is, if there is a linear order on \(\N\) invariant under every element of \(G\).
\begin{proposition}
    \label{prop:type-characterization}
    Let \(G\leq S_\infty\) be a closed subgroup. Then the following are equivalent.
    \begin{enumerate}
        \item[(i)] \(G\) is extremely amenable.
        \item[(ii)]
              \begin{enumerate}
                  \item[(a)] For every finite nonempty \(F\subseteq\N\),
                        \[
                            G_{(F)}=G_{\{F\}}.
                        \]
                  \item[(b)] Whenever \(\rho\leq \sigma\) are \(G\)-types and
                        \[
                            c\colon \rho\to\{1,\dots,k\}
                        \]
                        is a finite colouring, there exist \(i\in\{1,\dots,k\}\) and \(F\in \sigma\) such that
                        \[
                            c(F')=i
                            \qquad\text{for every } F'\subseteq F \text{ with } F'\in \rho.
                        \]
              \end{enumerate}
        \item[(iii)]
              \begin{enumerate}
                  \item[(a')] \(G\) preserves an ordering.
                  \item[(b)] The condition in (ii)(b) holds.\end{enumerate}
    \end{enumerate}
\end{proposition}

\begin{proof}


    Assume first that \(G\) is extremely amenable. Then \(G\) preserves an ordering because \(\LO(\N)\) is a compact \(G\)-flow. This gives (iii)(a').

    To prove (iii)(b), let \(\rho\leq \sigma\) be \(G\)-types and let \(c\colon \rho\to\{1,\dots,k\}\) be a colouring.
    Choose \(F'\in\rho\) and let \(V=G_{(F')}\).
    Since \(G\) preserves an ordering, every finite setwise stabiliser equals the corresponding pointwise stabiliser, so \(G_{(F')}=G_{\{F'\}}\).
    Hence \(G/V \cong G\cdot F'=\rho\).
    Now fix \(F_\sigma\in\sigma\) and consider the finite subset \(A=\{D\subseteq F_\sigma : D\in\rho\}\subseteq \rho\).
    Applying Proposition~\ref{prop:coset-characterization}, we obtain \(g\in G\) and \(i\in\{1,\dots,k\}\) such that \(c(g\cdot D)=i\) for all \(D\in A\).
    Let \(F=g\cdot F_\sigma\in \sigma\).
    If \(F'\subseteq F\) and \(F'\in\rho\), then \(g^{-1}\cdot F'\subseteq F_\sigma\) and \(g^{-1}\cdot F'\in \rho\), so \(c(F')=i\).
    Thus (iii)(b) holds.

    The implication (iii) \(\Rightarrow\) (ii) is immediate, since an invariant linear order on \(\N\) forces every finite setwise stabiliser to equal the corresponding pointwise stabiliser.

    Finally, assume (ii).
    We verify condition (ii) of Proposition~\ref{prop:coset-characterization} for open subgroups of the form \(G_{(F)}\).
    By (ii)(a), these equal the setwise stabilisers \(G_{\{F\}}\).
    Fix a colouring \(c\colon G/G_{\{F\}} \to \{1,\dots,k\}\) and a finite subset \(A\subseteq G/G_{\{F\}}\).
    Identify \(G/G_{\{F\}} \cong G\cdot F = \rho\).
    Let \(F_\sigma=\bigcup A\) and let \(\sigma\) be the \(G\)-type of \(F_\sigma\).
    Then \(\rho\leq \sigma\).
    By (ii)(b), there exist \(i\in\{1,\dots,k\}\) and \(E\in \sigma\) such that \(c(F')=i\) for every \(F'\subseteq E\) with \(F'\in \rho\).
    Since \(E\) and \(F_\sigma\) lie in the same \(G\)-type, we may write \(E=g\cdot F_\sigma\) for some \(g\in G\).
    Every element of \(A\) is a member of \(\rho\) contained in \(F_\sigma\), so after applying \(g\) it becomes a member of \(\rho\) contained in \(E\).
    Hence \(c(g\cdot a)=i\) for all \(a\in A\).
    By Proposition~\ref{prop:coset-characterization}, \(G\) is extremely amenable.
\end{proof}
We now isolate the corresponding Ramsey property for the group itself.

\begin{definition}
    \label{def:group-ramsey-property}
    Let \(G\leq S_\infty\). Suppose that \(\rho\leq \sigma\) are \(G\)-types. For \(F\in \sigma\), define
    \[
        \begin{pmatrix} F \\ \rho \end{pmatrix}
        =
        \{F'\subseteq F : F'\in \rho\}.
    \]
    If \(\rho\leq \sigma\leq \tau\) are \(G\)-types and \(k\geq 2\), we write
    \[
        \tau \to (\sigma)^\rho_k
    \]
    if for every \(F\in\tau\) and every colouring
    \(c\colon \begin{pmatrix} F \\ \rho \end{pmatrix}\to\{1,\dots,k\}\), there exists
    \(F_\sigma\in \begin{pmatrix} F \\ \sigma \end{pmatrix}\) such that
    \(c\restriction \begin{pmatrix} F_\sigma \\ \rho \end{pmatrix}\) is constant.

    We say that \(G\) has the \emph{Ramsey property} if for all \(G\)-types \(\rho\leq \sigma\) and all \(k\geq 2\), there exists a \(G\)-type \(\tau\geq \sigma\) such that
    \[
        \tau \to (\sigma)^\rho_k.
    \]
\end{definition}

\begin{theorem}
    \label{thm:ea-iff-ordering-plus-group-ramsey}
    Let \(G\leq S_\infty\) be a closed subgroup. Then the following are equivalent.
    \begin{enumerate}
        \item[(i)] \(G\) is extremely amenable.
        \item[(ii)] \(G\) preserves an ordering and has the Ramsey property.
    \end{enumerate}
\end{theorem}

\begin{proof}

    Assume that \(G\) is extremely amenable. Then \(G\) preserves an ordering by Proposition~\ref{prop:type-characterization}.

    We show that \(G\) has the Ramsey property. As in the usual Ramsey argument, it is enough to treat the case of two colours. Suppose instead that there exist \(G\)-types \(\rho\leq \sigma\) such that no \(G\)-type \(\tau\geq \sigma\) satisfies
    \[
        \tau\to (\sigma)^\rho_2.
    \]
    Fix \(F_\sigma\in \sigma\). Then for every finite set \(E\supseteq F_\sigma\) there exists a colouring \(c_E\colon \begin{pmatrix} E \\ \rho \end{pmatrix}\to\{1,2\}\) with no homogeneous member of \(\begin{pmatrix} E \\ \sigma \end{pmatrix}\).

    Let \(I\) be the set of all finite nonempty subsets of \(\N\), and choose an ultrafilter \(\mathcal U\) on \(I\) containing every set of the form \(\{E\in I : F\subseteq E\}\), where \(F\subseteq\N\) is finite. For each \(D\in \rho\), exactly one of the sets
    \[
        \{E\supseteq D\cup F_\sigma : c_E(D)=1\},
        \qquad
        \{E\supseteq D\cup F_\sigma : c_E(D)=2\}
    \]
    belongs to \(\mathcal U\).
    Define \(c(D)=i\) if the \(i\)th set belongs to \(\mathcal U\).
    This gives a colouring \(c\colon \rho\to\{1,2\}\).
    By Proposition~\ref{prop:type-characterization}, there exists \(F\in \sigma\) such that \(c\restriction \begin{pmatrix} F \\ \rho \end{pmatrix}\) is constant, say constantly equal to \(i\).
    For each \(D\in \begin{pmatrix} F \\ \rho \end{pmatrix}\), the set \(A_D=\{E\supseteq F\cup F_\sigma : c_E(D)=i\}\) belongs to \(\mathcal U\).
    Since there are only finitely many such \(D\), their intersection is nonempty.
    Choose
    \[
        E\in \bigcap_{D\in \begin{pmatrix} F \\ \rho \end{pmatrix}} A_D.
    \]
    Then \(E\supseteq F_\sigma\) and \(c_E(D)=i\) for all \(D\in \begin{pmatrix} F \\ \rho \end{pmatrix}\).
    Thus \(F\) is a homogeneous member of \(\begin{pmatrix} E \\ \sigma \end{pmatrix}\) for the colouring \(c_E\), contradicting the choice of \(c_E\).
    Therefore \(G\) has the Ramsey property.

    Conversely, assume that \(G\) preserves an ordering and has the Ramsey property. To prove extreme amenability, by Proposition~\ref{prop:type-characterization} it is enough to verify condition (ii)(b) there. Let \(\rho\leq \sigma\) be \(G\)-types and let \(c\colon \rho\to\{1,\dots,k\}\) be a colouring. Since \(G\) has the Ramsey property, there exists a \(G\)-type \(\tau\geq \sigma\) such that
    \[
        \tau\to (\sigma)^\rho_k.
    \]
    Choose \(F\in\tau\). Then the restriction of \(c\) to \(\begin{pmatrix} F \\ \rho \end{pmatrix}\) has a homogeneous member of \(\begin{pmatrix} F \\ \sigma \end{pmatrix}\), which is exactly the statement required in Proposition~\ref{prop:type-characterization}. Hence \(G\) is extremely amenable.
\end{proof}



\begin{remark}
    \label{rem:cofinal-types}
    A family \(T\) of \(G\)-types is called \emph{cofinal} if for every \(G\)-type \(\rho\) there exists \(\sigma\in T\) such that \(\rho\leq \sigma\).
    It is enough in Definition~\ref{def:group-ramsey-property} to consider \(\rho\), \(\sigma\), and \(\tau\) ranging over any fixed cofinal family of \(G\)-types.
\end{remark}



\section[Fraïssé order classes]{Fraïssé order classes and extremely amenable\\automorphism groups}
\label{sec:fraisse-order-classes-and-ea-groups}

We now translate the characterisation obtained in the previous section into Fraïssé-theoretic language.

Let \(L\) be a signature containing a distinguished binary relation symbol \(<\). An \emph{order structure} for \(L\) is an \(L\)-structure \(\bfa\) such that its interpretation \(\prec^{\bfa}\) of \(<\) is a linear ordering on \(A\). A class \(\cK\) of finite \(L\)-structures is called an \emph{order class} if every member of \(\cK\) is an order structure.

The next theorem identifies the extremely amenable closed subgroups of \(S_\infty\) with automorphism groups of Fraïssé limits of Ramsey order classes.

\begin{theorem}
    \label{thm:ea-iff-fraisse-order-ramsey}
    Let \(G\leq S_\infty\) be a closed subgroup. Then the following are equivalent.
    \begin{enumerate}
        \item[(i)] \(G\) is extremely amenable.
        \item[(ii)] There is a Fraïssé order class \(\cK\) with the Ramsey property such that \(G=\Aut(\bff)\), where \(\bff=\Flim(\cK)\).
    \end{enumerate}
\end{theorem}

\begin{proof}


    Assume first that \(G\) is extremely amenable. Let \(\bfa_G\) be the canonical structure on \(\N\) associated with \(G\); recall from Chapter~\ref{chap:fraisse-automorphism-groups} that its basic \(n\)-ary relations are the \(G\)-orbits on \(\N^n\), and that \(\Aut(\bfa_G)=G\).
    By construction, the structure \(\bfa_G\) is ultrahomogeneous and has a relational signature.

    Since \(G\) is extremely amenable, Theorem~\ref{thm:ea-iff-ordering-plus-group-ramsey} shows that \(G\) preserves an ordering on \(\N\).
    Let \(\prec^{\bfa}\) be a linear ordering of \(\N\) fixed by every element of \(G\).
    Expand \(\bfa_G\) by a new binary relation symbol \(<\) interpreted as \(\prec^{\bfa}\), and denote the resulting structure by \(\bfa=\langle \bfa_G,\prec^{\bfa}\rangle\).
    Then every element of \(G\) preserves \(\prec^{\bfa}\), so \(G\leq \Aut(\bfa)\).
    Conversely, every automorphism of \(\bfa\) is in particular an automorphism of \(\bfa_G\), hence belongs to \(G\).
    Therefore \(\Aut(\bfa)=G\).

    Since \(\bfa_G\) is ultrahomogeneous and \(\prec^{\bfa}\) is \(G\)-invariant, the expansion \(\bfa\) is also ultrahomogeneous: if \(\varphi\colon \bfb\hookrightarrow \bfc\) is an isomorphism between finite substructures of \(\bfa\), then \(\varphi\) is in particular an isomorphism between finite substructures of \(\bfa_G\), so it extends to some \(g\in G=\Aut(\bfa_G)\); because \(g\) preserves \(\prec^{\bfa}\), it is an automorphism of \(\bfa\) extending \(\varphi\).
    Because the signature is relational, finitely generated substructures of \(\bfa\) are finite, so \(\bfa\) is locally finite.
    It follows that \(\cK:=\Age(\bfa)\) is a Fraïssé order class with Fraïssé limit \(\bfa\).

    It remains to show that \(\cK\) has the Ramsey property.
    Let \(\bfb,\bfc\in\cK\) with \(\bfb\hookrightarrow \bfc\), and let \(k\geq 2\).
    Consider the \(G\)-types \(\rho_{\bfb}\) and \(\rho_{\bfc}\) consisting of all subsets of \(\N\) which, with the induced structure from \(\bfa\), are isomorphic to \(\bfb\) and \(\bfc\), respectively.
    Since \(\bfa\) is ultrahomogeneous, copies of a finite substructure in \(\bfa\) form a single \(G\)-type.
    By Theorem~\ref{thm:ea-iff-ordering-plus-group-ramsey}, the group \(G\) has the Ramsey property.
    Hence there exists a \(G\)-type \(\tau\geq \rho_{\bfc}\) such that \(\tau\to (\rho_{\bfc})^{\rho_{\bfb}}_k\).
    Choose a finite substructure \(\bfd\leq \bfa\) whose underlying set belongs to \(\tau\).
    Then \(\bfd\in\cK\), and the statement \(\bfd\to (\bfc)^{\bfb}_k\) is exactly the same as \(\tau\to (\rho_{\bfc})^{\rho_{\bfb}}_k\).
    Thus \(\cK\) has the Ramsey property.

    This proves (i) \(\Rightarrow\) (ii).

    Now assume (ii). Let \(\bff=\Flim(\cK)\) and \(G=\Aut(\bff)\), where \(\cK\) is a Fraïssé order class with the Ramsey property.

    Since \(\bff\) is an order structure, its linear order \(\prec^{\bff}\) on \(F\) is preserved by every automorphism of \(\bff\). Thus \(G\) preserves an ordering.

    To show that \(G\) is extremely amenable, it is enough by Theorem~\ref{thm:ea-iff-ordering-plus-group-ramsey} to verify that \(G\) has the Ramsey property. By Remark~\ref{rem:cofinal-types}, it suffices to check this on any cofinal family of \(G\)-types. We take the family of \(G\)-types arising from finite substructures of \(\bff\).

    Let \(\bfb,\bfc\in\Age(\bff)=\cK\) with \(\bfb\hookrightarrow \bfc\), and let \(\rho_{\bfb},\rho_{\bfc}\) be the corresponding \(G\)-types of copies of \(\bfb\) and \(\bfc\) in \(\bff\). Since \(\cK\) has the Ramsey property, for every \(k\geq 2\) there exists \(\bfd\in\cK\) such that \(\bfd\to (\bfc)^{\bfb}_k\). Let \(\tau\) be the \(G\)-type of copies of \(\bfd\) in \(\bff\). Then \(\tau\to (\rho_{\bfc})^{\rho_{\bfb}}_k\).
    Therefore \(G\) has the Ramsey property on this cofinal family of types, hence \(G\) has the Ramsey property. Since \(G\) also preserves an ordering, Theorem~\ref{thm:ea-iff-ordering-plus-group-ramsey} yields that \(G\) is extremely amenable.
\end{proof}

The preceding theorem immediately gives the Fraïssé-theoretic form that will be used in the next chapter.

\begin{corollary}
    \label{cor:fraisse-order-class-ea}
    Let \(\cK\) be a Fraïssé order class and let \(\bff=\Flim(\cK)\).
    Then the following are equivalent.
    \begin{enumerate}
        \item[(i)] \(\Aut(\bff)\) is extremely amenable.
        \item[(ii)] \(\cK\) has the Ramsey property.
    \end{enumerate}
\end{corollary}

\begin{proof}
    Apply Theorem~\ref{thm:ea-iff-fraisse-order-ramsey} with \(G=\Aut(\bff)\).
\end{proof}

The preceding corollary is the form of the Kechris--Pestov--Todorcevic correspondence that will be used in the next chapter. It shows that, for a Fraïssé order class, the extreme amenability of the automorphism group of its Fraïssé limit is exactly equivalent to the Ramsey property of the class. In Chapter~\ref{chap:examples-extreme-amenability} we apply this criterion to concrete examples.
